\chapter{A novel analysis on learning theory}
\section{Introduction}
Motivated by earlier analysis and observations, it is apparent that the analysis is entangled, with various situations, interpretations, setting and assumptions together, and with a non-unified structure of models in question. The system in which the model is considered and analyzed is also ineffective when considering larger structures or more `realistic' consideration aside from very simplified, strong system itself. There are a lot of ambiguities around what is considered and what is not, even with probabilistic nature and the deterministic aspect. Terms like \textbf{latent spaces} and many do not capture the entire dynamic or meaning of the system. It is then we introduce our novel analytical framework, namely the \textbf{Unified Unitary Construct} framework. It is however, can be said otherwise that the learning theory is rather not so novel, but approaching the problem is a more distinct way of the already existing framework. 

\subsection{Principle}

We are motivated from earlier analysis, further experimental results, and arguably practical observations of empirical-based architectures that we find interesting. Oddly enough, the idea is encouraged of the observation that \textit{deep neural network} organising architecture might be more general than it is, particularly in the sense of mathematical modelling and structural implicit approximation, as well as factors like scalability, generality, and else. Furthermore, one of the purpose of this novel theory is also to redefine the entire learning space to, at least if not entirely original, might as well fall into the category of a reformulation of theoretical framework and not partaking the role of replacing it. 

As such, we need some principles to get started the analysis. What can those be that guarantee the success of a particular reformulation at first? For starter, we indict on the requirement of developing a general framework of neural network-like organisation for our models, both in the interpretation of the learning objective, and the structure of the hypothesis learner. Secondly, we have to define learning as a barebone process, of which then we will have to justify \textit{probabilistic nature} of it and the degree in which that is true. Thereby, we will then also have to reconsider the learning problem, under what scope of the problem and the classes of learning that is there to be achieved, the encoding space and what it is meant to, and the criteria of guaranteeing such learning to be probable. For now, those are enough, but the development that come afterward will justify it instead of the principle. 
Nevertheless, we can see the principle here - we want to redefine learning process in an organized manner, with minimal messiness from considering all factors together, and normalize the interpretation of \textit{unit-wise} architecture that neural network uses up to more. 

\subsection{Planning}

Then, what is the plan? The goal is to \textit{systematically introduce} the concept and the novel theory. Hence, we will separate it into discrete small step size parts. First, we need to conform the setting with a new template. This mean the following order is needed. 
\begin{enumerate}[leftmargin=1cm,itemsep=2pt,topsep=2pt]
  \item Introduce and define the formalism of an \textit{ambiance space} and an \textit{encoding space}. Those two will come in together to form the majority of $S$ as per system dynamic notation. 
  \item In the \textit{collection} $S$ should we extend it further, there always exists the object $h$, belonging to certain configuration class $\mathcal{H}$. We need to define this particular object, and its interaction to the collection itself. Which means, the ambiance space and encoding space. How do we bring up the \textit{representation concept} here altogether? 
  \item To start the mechanics, we need the \textit{mechanical space} $M$ to be defined. What can be it? What connects and form the structural description that $S$ make with the ambiance space, encoding, representation, subject dynamics, and else? 
  \item Now, come to the \textit{operational process}. For this, let's assume we have two different systems, in accordances, $S_{1}$ and $S_{2}$, each of which houses $h$, and $c$ respectively. How do they communicate? How, should it be done, can the operational space indict particular objective, or requirement, such that $h$ is evaluated, and $c$ is evaluated? What is hidden of $h$, and what is hidden of $c$? 
  \item Then, finally, bounds, theorems, analytical results, algorithmic consideration and conjectures that fit, into this description. 
\end{enumerate}
From the scope of it, it might seem not too practical to call as novel. But as far as we are considered, there should be something to take account on, and we need to do them in a careful manner, such is to say there has been none, of which partake in said direction. There is a reason why messy construction being the worst nightmare of any people using the house, and we are the one to fix it, or at least attempted so. 